% Source: https://artofproblemsolving.com/community/c3h1706974_simplifying_in_modular_arithmetic
% Stable identity: aops-topic-1706974-q01
A congruence can be treated like an equation as far as adding or subtracting the same integer on both sides is concerned: if $x\equiv y\pmod n$ then $x-k\equiv y-k\pmod n$, since $(x-k)-(y-k)=x-y$ is still divisible by $n$.

\textbf{(a)} Subtract $4$ from both sides: $b\equiv9-4=5\pmod{20}$. So the solutions are $\boxed{b\equiv5\pmod{20}}$, i.e.\ $b=\dots,-15,5,25,45,\dots$. The student's $b\equiv4$ is wrong: for $b=4$, $b+4=8\not\equiv9\pmod{20}$.

\textbf{(b)} Subtract $2$ from both sides: $b\equiv1-2=-1\pmod3$. Since $-1\equiv2\pmod3$ (they differ by $3$), this is $\boxed{b\equiv2\pmod3}$. Check: $b=2$ gives $b+2=4\equiv1\pmod3$.
